\subsection{The Krein-Rutman theorem}

Lemma 3. --- Let \((a_n)_{n \geqslant 0}\) be a sequence of positive real numbers such that the power series

\[
f (z) = \sum_ {n \geqslant 0} a _ {n} z ^ {n}
\]

has a finite radius of convergence r \textgreater{} 0. Let D ⊂ C be the open disk with center 0 and radius r. There does not exist a holomorphic function \(\widetilde{f}\) defined on an open neighborhood U of r in C that coincides with f on U ∩ D.

Suppose that such a holomorphic function \(\tilde{f}\), defined on an open neighborhood U of \(r\), exists. There exist real numbers \(s < r\) and \(\delta > 0\) such that \(s + \delta > r\) and such that the open disk D' with center \(s\) and radius \(\delta\) is contained in U. The power series expansion of \(\tilde{f}\) at the point \(s\) (VAR, R1, p. 26, 3.2.1) then converges for every \(z\) in the disk with center 0 and radius \(\delta\) (VAR, R1, p. 29, 3.3.4). This series is

\[
\widetilde {f} _ {s} (z) = \sum_ {n = 0} ^ {+ \infty} \frac {\widetilde {f} ^ {(n)} (s)}{n !} z ^ {n} = \sum_ {n = 0} ^ {+ \infty} \frac {f ^ {(n)} (s)}{n !} z ^ {n}
\]

(VAR, R1, p. 27, 3.2.4). Since \(s \in D\) , we have

\[
f ^ {(n)} (s) = \sum_ {k = n} ^ {+ \infty} k (k - 1) \dots (k - n + 1) a _ {k} s ^ {k - n}
\]

(VAR, R1, p. 28, 3.2.11). Take \(z\) such that \(0 < z < \delta\) . Since \(a_{k} \geqslant 0\) , we have

\[
\begin{array}{r l} & {\widetilde {f} _ {s} (z) = \sum_ {n = 0} ^ {+ \infty} \biggl (\sum_ {k = n} ^ {+ \infty} \frac {k (k - 1) \cdots (k - n + 1)}{n !} a _ {k} s ^ {k - n} \biggr) z ^ {n}} \\ & {\qquad = \sum_ {k = 0} ^ {+ \infty} \biggl (a _ {k} \sum_ {n = 0} ^ {k} \frac {k !}{n ! (k - n) !} s ^ {k - n} z ^ {n} \biggr) = \sum_ {k = 0} ^ {+ \infty} a _ {k} (s + z) ^ {k}} \end{array}
\]

(TG, III, p. 40). When \(s + z > r\) , the convergence of this expression contradicts the hypothesis that the radius of convergence of the power series \(f\) is equal to \(r\) .

Let E be a real Banach space. Let C be a pointed convex cone in E. The vector space generated by C is equal to C -- C (EVT, II, p. 12, cor. 1). In particular, the cone C is total in E if and only if C -- C is dense in E. Recall that C is said to be salient if C ∩ (−C) is reduced to 0 (EVT, II, p. 11).

The polar C° of C is the set of continuous linear forms \(\ell \in E'\) such that \(\ell(x) \geqslant 0\) for all \(x \in C\) (EVT, II, p. 47, prop. 4). By the bipolar theorem (EVT, II, p. 48, th. 1), if C is a closed pointed convex cone, then C°° = C.

Lemma 4. --- Let E be a real Banach space and \(u \in \mathcal{L}(\mathrm{E}_{(\mathbf{C})})\) a nonzero endomorphism of \(\mathrm{E}_{(\mathbf{C})}\) . Let C be a total convex cone in E. There exists \(x \in \mathrm{C}\) such that \(u(x) \neq 0\) .

Indeed, if \(u\) vanishes on C, then \(u\) vanishes on C - C, hence on E, and therefore on \(\mathrm{E}_{(\mathbf{C})}\) .

THEOREM 4 (Krein--Rutman). --- Let E be a complete normable space over R. Let C be a closed total salient convex cone in E and let \(u \in \mathcal{L}(\mathrm{E})\) be a compact linear map such that \(u(\mathrm{C}) \subset \mathrm{C}\) . If the spectral radius \(\varrho(u)\) is \textgreater0, then \(\varrho(u)\) is an isolated point of the spectrum of u, and there exists a nonzero eigenvector \(x \in C\) of u for the eigenvalue \(\varrho(u)\) .

Let \(E_{(\mathbf{C})}\) be the complexification of E and \(u_{(\mathbf{C})}\) the endomorphism of \(E_{(\mathbf{C})}\) obtained by extension of scalars from u. The spectral radius of \(u_{(\mathbf{C})}\) is equal to \(\varrho(u)\) (I, p. 86); it is simply denoted by \(\varrho\) . Since \(\varrho > 0\) , the complex spectrum of \(u\) is not reduced to 0. There therefore exists \(\lambda_0 \in \mathrm{Sp}_\mathrm{s}(u_{(\mathbf{C})})\) such that \(|\lambda_0| = \varrho\) (prop. 5 of III, p. 90). Let \(e_0\) be the spectral projector of \(u_{(\mathbf{C})}\) associated with \(\lambda_0\) .

The resolvent \(\lambda \mapsto \mathrm{R}(u_{(\mathbf{C})},\lambda) = (\lambda -u_{(\mathbf{C})})^{-1}\) is holomorphic on the complement of the spectrum of \(u_{(\mathbf{C})}\) (th. 1 of I, p. 24). The complex number \(\lambda_0\) is a pole of the resolvent and its residue is the spectral projector \(e_0\) (I, p. 131).

Let \(y\) and \(y'\) be elements of E such that \(y + iy' \in \mathrm{E}_{(\mathbf{C})}\) is an eigenvector of \(u_{(\mathbf{C})}\) for the eigenvalue \(\lambda_0\) . Since \(e_0(y + iy') = y + iy' \neq 0\) , there exists an element \(x_0 \in \mathrm{C}\) such that \(e_0(x_0) \neq 0\) (lemma 4). Since C is closed and pointed, its polar \(\mathrm{C}^\circ\) is total (EVT, II, p. 48, cor. 1), and hence there exists a linear form \(\ell_0 \in \mathrm{C}^\circ\) such that \(\langle e_0(x_0), \ell_0 \rangle \neq 0\) .

Consider the function \(f\) defined by \(f(0)=0\) and by

\[
f (z) = \langle \mathrm{R} (u _ {(\mathbf {C})}, z ^ {- 1}) x _ {0}, \ell_ {0} \rangle
\]

for \(z \in \mathbf{C}\) such that \(z^{-1} \notin \mathrm{Sp}(u_{(\mathbf{C})})\) . This function satisfies

\[
f (z) = \ell_ {0} \bigg (\Bigl (\sum_ {n = 0} ^ {+ \infty} z ^ {n + 1} u _ {(\mathbf {C})} ^ {n} \Bigr) x _ {0} \bigg) = \sum_ {n = 0} ^ {\infty} \langle u ^ {n} (x _ {0}), \ell_ {0} \rangle z ^ {n + 1}\tag{1}
\]

for \(|z| < 1 / \varrho\) (Theorem 1 of I, p. 24, d)) and is therefore holomorphic in the disk centered at 0 with radius \(1 / \varrho\) .

There exists a holomorphic function \(\widetilde{\mathrm{R}}\) defined on an open neighborhood U of \(\lambda_0\) and with values in \(\mathcal{L}(\mathrm{E})\) such that for \(z\) belonging to \(\mathrm{U} - \{\lambda_0\}\) , we have

\[
\mathrm{R} (u _ {(\mathbf {C})}, z) = \widetilde {\mathrm{R}} (z) + \sum_ {n = 0} ^ {+ \infty} (z - \lambda_ {0}) ^ {- n - 1} (u _ {(\mathbf {C})} - \lambda_ {0}) ^ {n} e _ {0}
\]

(prop. 17 of I, p. 83). For \(z\) such that \(z^{-1} \in \mathrm{U}\) and \(z \neq 1 / \lambda_0\) , we therefore have

\[
f (z) = \langle \widetilde {\mathrm{R}} (z ^ {- 1}) x _ {0}, \ell_ {0} \rangle + \sum_ {n = 0} ^ {+ \infty} (z ^ {- 1} - \lambda_ {0}) ^ {- n - 1} \langle (u _ {(\mathbf {C})} - \lambda_ {0}) ^ {n} e _ {0} (x _ {0}), \ell_ {0} \rangle .
\]

The term of the series corresponding to \(n=0\) is \((z^{-1}-\lambda_{0})^{-1}\langle e_{0}(x_{0}),\ell_{0}\rangle\) . Since \(\langle e_{0}(x_{0}),\ell_{0}\rangle\neq0\) , the uniqueness of the Laurent expansion (VAR, R1, p. 30, 3.3.9) implies that \(f\) does not extend to a holomorphic function in a neighborhood of \(1/\lambda_{0}\) . In particular, the radius of convergence of the power series expansion (1) of the function \(f\) at the point 0 is equal to \(1/\varrho\) .

The coefficients of the power series (1) are \(\langle u^n(x_0),\ell_0\rangle\geqslant 0\) since \(u(\mathrm{C})\subset\mathrm{C}\) and \(\ell_{0}\in\mathrm{C}^{\circ}\) . By Lemma 3, the function \(f\) does not extend to a holomorphic function in a neighborhood of \(1/\varrho\) . The resolvent of \(u_{(\mathbf{C})}\) therefore cannot be extended to a holomorphic function in a neighborhood of \(\varrho\) , that is, \(\varrho\in\mathrm{Sp}(u)\) . This implies that \(\varrho\) is an eigenvalue of \(u_{(\mathbf{C})}\) (prop. 5 of III, p. 90). Since \(\varrho\) is real, it is also an eigenvalue of \(u\) .

Let \(d \geqslant 1\) be the order of the pole of the resolvent of \(u_{(\mathbf{C})}\) at \(\varrho\) . Let \(e\) be the endomorphism

\[
e = \lim _ {z \to \varrho} (z - \varrho) ^ {d} \mathrm{R} (u _ {(\mathbf {C})}, z).
\]

It is nonzero and commutes with \(u_{(\mathbf{C})}\) , and its image is contained in the eigenspace of \(u_{(\mathbf{C})}\) relative to \(\varrho\) (cf. no. 3 of I, p. 131). Now let \(\ell\) be an element of C° and \(x\) an element of C. By Theorem 1 of I, p. 24, d), we have

\[
\langle e(x),\ell \rangle = \lim_{\substack{z\to \varrho \\ z > \varrho}}(z - \varrho)^{d}\sum_{n\geqslant 0}\langle u^{n}(x),\ell \rangle z^{-n - 1}\geqslant 0.
\]

The bipolar theorem (EVT, II, p. 48, th. 1) implies that \(e(x) \in \mathrm{C}\) for all \(x \in \mathrm{C}\) . Since \(\mathrm{C}\) is total, there exists \(x \in \mathrm{C}\) such that \(e(x) \neq 0\) (Lemma 4), and then \(e(x)\) belongs to \(\mathrm{C}\) and is an eigenvector of \(u\) for the eigenvalue \(\varrho\) , as desired.

COROLLARY. --- Let E be a complete normable space over R. Let C be a closed, total, salient convex cone in E, and \(u \in \mathcal{L}(E)\) be a compact linear map such that \(u(\mathrm{C}) \subset \mathrm{C}\) . If the spectral radius \(\varrho(u)\) of u is \textgreater{} 0, then \(\varrho(^{t}u) = \varrho(u)\) is an eigenvalue of \(^{t}u\) , and \(^{t}u\) admits an eigenvector relative to \(\varrho(u)\) in \(C^{\circ}\) .

One has \(\varrho(^{t}u) = \varrho(u)\) (prop. 3 of I, p. 131). By corollary 1 of III, p. 9, the endomorphism \(^{t}u\) of E' is compact. Moreover, we have \(^{t}u(C^{\circ}) \subset C^{\circ}\) ; the assertion then follows from the Krein-Rutman theorem applied to \(^{t}u\) , and to the closed convex cone \(C^{\circ}\) since the latter is salient (because C is total) and total (because C is salient).

Remark. --- The hypothesis \(\varrho(u) > 0\) cannot be omitted in general in the Krein–Rutman theorem. For example, let V be the endomorphism of the Banach space \(\mathcal{C}([0,1])\) defined by

\[
\mathrm{V} (f) (x) = \int_ {0} ^ {x} f (y) d y
\]

for \(f \in \mathcal{C}([0,1])\) . The map V is compact and its spectral radius is zero (exercise 1 of I, p. 187). It preserves the closed, total, salient convex cone in \(\mathcal{C}([0,1])\) consisting of positive functions, and has no eigenvalue (loc. cit.).

The following proposition describes a sufficient condition for an endomorphism preserving a cone to have a strictly positive spectral radius, and it then refines the Krein-Rutman theorem.

PROPOSITION 8. --- Let E be a nonzero complete normable space over R. Let C be a closed salient convex cone with nonempty interior \(\mathring{\mathrm{C}}\) in E, and let \(u \in \mathcal{L}(\mathrm{E})\) be a compact linear map such that \(u(\mathrm{C} - \{0\}) \subset \mathring{\mathrm{C}}\) .

\begin{enumerate}
\def\labelenumi{\alph{enumi})}
\item
  One has \(\varrho(u) > 0\) and there exists an eigenvector \(x_0\) of \(u\) in \(\mathring{C}\) for the eigenvalue \(\varrho(u)\) ;
\item
  The spectral projector of u corresponding to \(\varrho(u)\) has rank 1;
\item
  For every eigenvalue \(\lambda \neq \varrho(u)\) of \(u_{(\mathbf{C})}\) , one has \(|\lambda| < \varrho(u)\) ;
\item
  Let F be a subspace of E stable under u such that \((\mathrm{C}-\{0\})\cap\mathrm{F}\) is nonempty. Then \(x_{0}\in F\) . In particular, the only eigenvectors of u in C are the multiples of \(x_{0}\) .
\end{enumerate}

Let us prove two preliminary lemmas.

Lemma 5. --- Let E be a real Banach space, C a convex cone in E, and \(u \in \mathcal{L}(\mathrm{E})\) such that \(u(\mathrm{C} - \{0\}) \subset \mathring{\mathrm{C}}\) . Let \(\ell \in \mathrm{C}^{\circ}\) be an eigenvector of \(^{t}u\) . Then the kernel of \(\ell\) is stable under \(u\) and does not meet \(\mathrm{C} - \{0\}\) .

Let \(\lambda \in \mathbf{R}\) such that \(^t u(\ell) = \lambda \ell\) . For every \(x \in \mathrm{Ker}(\ell)\) , we have

\[
\langle u (x), \ell \rangle = \langle x, ^ {t} u (\ell) \rangle = \lambda \langle x, \ell \rangle = 0,
\]

hence \(\operatorname{Ker}(\ell)\) is stable under \(u\) .

Let \(x\) be a nonzero element of \(\operatorname{Ker}(\ell)\) . For every element \(y\) of E such that \(\langle y, \ell \rangle < 0\) , we have \(\langle u(x) + y, \ell \rangle < 0\) , whence it follows that \(u(x) + y \notin \mathrm{C}\) since \(\ell \in \mathrm{C}^\circ\) . We deduce that \(u(x)\) does not belong to \(\mathring{\mathrm{C}}\) . Since \(u(\mathrm{C} - \{0\}) \subset \mathring{\mathrm{C}}\) , this implies that \(x \notin \mathrm{C}\) .

Lemma 6. --- Let E be a finite-dimensional real vector space and let B be a compact convex neighborhood of 0 in E. Let u be an endomorphism of E such that \(u(B) \subset \mathring{B}\) . The complex spectrum of u is then contained in the unit disk of C and, in particular, it does not meet the circle with center 0 and radius 1 in C.

By replacing B with \(B \cap (-B)\) , we reduce to the case where B is balanced. The gauge p of B (EVT, II, p. 28, ex. 3) is then a norm on E, which defines its topology (EVT, I, p. 14, th. 2). The hypothesis implies that \(p(u(x)) < p(x)\) for every element x of E - \{0\}, hence that the spectral radius of u is \textless{} 1; since this is also the spectral radius of \(u_{(\mathbf{C})}\) (cf. I, p. 86), the conclusion follows.

Let us now prove the proposition.

Denote by \(x \preccurlyeq y\) the order relation on E associated with the convex cone C (EVT, II, p. 13, n° 5), that is, \(x \preccurlyeq y\) if and only if \(y - x \in \mathrm{C}\) . We have \(u(x) \preccurlyeq u(y)\) if \(x \preccurlyeq y\) . Since \(\mathring{\mathrm{C}}\) is nonempty, the vector space C -- C generated by C (EVT, II, p. 12, cor. 1) contains a neighborhood of 0, hence the cone C is total.

Let us prove assertion \(a\) ). Let \(\varrho = \varrho(u)\) . Let \(y_0\) be an element of \(\mathring{\mathrm{C}}\) . We have \(y_0 \neq 0\) . Let \(r > 0\) be such that the closed ball with center \(y_0\) and radius \(r\) is contained in C. For every \(y \in \mathrm{E} - \{0\}\) , we have \(y_0 - r \| y\|^{-1}y \in \mathrm{C}\) , hence \(y \preccurlyeq r^{-1}\| y\| y_0\) for all \(y \in \mathrm{E}\) .

Since \(y_{0} \neq 0\) , the hypotheses imply that \(u(y_{0}) \in \mathring{\mathrm{C}}\) . There therefore exists t \textgreater{} 0 such that \(tu(y_{0}) - y_{0} \in \mathrm{C}\) . Set v = tu. This is a compact endomorphism of E such that \(v(\mathrm{C}) \subset \mathrm{C}\) and \(v(y_{0}) \succcurlyeq y_{0}\) . For every integer \(n \geqslant 1\) , we have

\[
y _ {0} \preccurlyeq v ^ {n} (y _ {0}) \preccurlyeq r ^ {- 1} \| v ^ {n} (y _ {0}) \| y _ {0} \preccurlyeq r ^ {- 1} \| v ^ {n} \| \| y _ {0} \| y _ {0},
\]

hence \((r^{-1}\|v^{n}\|\|y_{0}\|-1)y_{0}\in\mathrm{C}\) . Since C is salient, this implies that \(t^{n}\|u^{n}\|=\|v^{n}\|\geqslant r/\|y_{0}\|\) , whence \(\varrho\geqslant t^{-1}>0\) (I, p. 20, prop. 1).

By the Krein--Rutman theorem (th. 4), the real number \(\varrho\) is an eigenvalue of \(u\) , and there exists an eigenvector \(x_0\) of \(u\) in C for the eigenvalue \(\varrho\) . We have \(\varrho x_0 = u(x_0) \in \mathring{\mathrm{C}}\) by hypothesis, hence \(x_0 \in \mathring{\mathrm{C}}\) . This establishes assertion \(a\) ).

Let K be the intersection of the spectrum of \(u_{(\mathbf{C})}\) and the circle with center 0 and radius \(\varrho\) in C. Since u is compact and \(\varrho > 0\) , the set K is finite and the image of the spectral projector \(e_{K}\) is a finite-dimensional subspace of \(\mathrm{E}_{(\mathbf{C})}\) (prop. 2 of III, p. 83 and prop. 5 of III, p. 90). Since K is stable under complex conjugation, the image of \(e_{K}\) is a finite-dimensional subspace of \(\mathrm{E}_{(\mathbf{C})}\) rational over R (A, V, p.60, prop. 6); denote by F the subspace of E such that \(\mathrm{F}_{(\mathbf{C})}\) is equal to the image of \(e_{K}\) . The space F is stable under u and contains the eigenspace of u relative to \(\varrho\) , hence F is nonzero.

To prove the assertions \(b)\) and \(c),\) it suffices to prove that \(F\) is of dimension 1.

Denote by \(v\) the endomorphism of F deduced from \(u\) by passing to subspaces. Let \(C_F = C \cap F\) . This is a closed salient convex cone with nonempty interior in F (since \(x_0 \in \mathring{C} \cap F\) ), hence total; it satisfies \(v(C_F - \{0\}) \subset \mathring{C}_F\) , and in particular is stable under \(v\) . Since \(\varrho(v) \leqslant \varrho\) and \(x_0 \in F\) , we have \(\varrho(v) = \varrho\) . By the corollary of Theorem 4, applied to \(C_F\) and \(v\) there exists an eigenvector \(\ell\) of \(^t v\) in \(C_F^\circ\) for the eigenvalue \(\varrho(v) = \varrho > 0\) .

The subspace \(\mathrm{Ker}(\ell)\) of F is stable under v. Let w be the endomorphism of \(\mathrm{Ker}(\ell)\) deduced from \(\varrho^{-1}v\) by passing to the quotient subspaces; its spectrum is contained in the unit circle. The set B of \(y \in \mathrm{Ker}(\ell)\) such that \(x_{0} + y \in C_{F}\) is a closed convex neighborhood of 0 in \(\mathrm{Ker}(\ell)\) . For every \(y \in B\) and every \(z \in \mathrm{Ker}(\ell)\) , we have

\[
x _ {0} + (w (y) + z) = \varrho^ {- 1} (v (x _ {0} + y) + \varrho z)
\]

which belongs to \(C_{F}\) if the norm of z is sufficiently small, hence \(w(y) \in \mathring{\mathrm{B}}\) . The set B is bounded: indeed, if there existed a sequence \((y_{n})_{n \in \mathbf{N}}\) in \(\operatorname{Ker}(\ell)\) such that \(\|y_{n}\| \to +\infty\) and \(x_{0} + y_{n} \in C_{F}\) then, after passing to a subsequence, we would have \(y_{n}/\|y_{n}\| \to y\) where \(y \in \operatorname{Ker}(\ell)\) is nonzero, and \(y = \lim \|y_{n}\|^{-1}(x_{0} + y_{n})\) would belong to \(C_{F}\) , contradicting Lemma 5. The set B is therefore compact.

One then deduces from Lemma 6, applied to \(\mathrm{Ker}(\ell)\) , to B, and to \(w\) , that the complex spectrum of \(w\) does not meet the circle with center 0 and radius 1; this implies that the spectrum of \(w\) is empty, which means that \(\mathrm{Ker}(\ell)\) is reduced to \(\{0\}\) , hence that F has dimension 1. The assertions \(b\) and \(c\) are therefore established.

The linear map \({}^t u\) is compact (Corollary 1 of III, p. 9) and \({}^t u(\mathrm{C}^\circ) \subset \mathrm{C}^\circ\) ; moreover \(\varrho(^t u) = \varrho > 0\) (Prop. 3 of I, p. 131). By the corollary of Theorem 4, there exists in \(\mathrm{C}^\circ\) an eigenvector \(\ell \neq 0\) of \({}^t u\) for the eigenvalue \(\varrho\) . The kernel of \(\ell\) is stable under \(u\) and does not meet \(\mathrm{C} - \{0\}\) (Lemma 5). In particular, one has \(x_0 \notin \operatorname{Ker}(\ell)\) .

Let F be a subspace of E stable under u. One has the decomposition \(\mathrm{F} = (\mathrm{F} \cap \mathbf{R}x_{0}) \oplus (\mathrm{F} \cap \operatorname{Ker}(\ell))\) . If F does not contain \(x_{0}\) , we therefore have \(\mathrm{F} \subset \operatorname{Ker}(\ell)\) , then F does not meet \(C - \{0\}\) . This establishes d) and concludes the proof of the proposition.

COROLLARY. --- Let E be a nonzero complete normable space over R. Let C be a closed convex salient cone with nonempty interior \(\dot{C}\) in E, and let u be a compact endomorphism of E such that \(u(\mathrm{C}) \subset \mathrm{C}\) . Suppose that there exists an integer \(k \geqslant 1\) such that \(u^{k}(\mathrm{C} - \{0\}) \subset \mathring{\mathrm{C}}\) .

\begin{enumerate}
\def\labelenumi{\alph{enumi})}
\item
  One has \(\varrho(u) > 0\) and there exists an eigenvector \(x_0\) of \(u\) in \(\dot{C}\) for the eigenvalue \(\varrho(u)\) ;
\item
  The spectral projector of u corresponding to \(\varrho(u)\) has rank 1;
\item
  For every eigenvalue \(\lambda \neq \varrho(u)\) of \(u_{(\mathbf{C})}\) , one has \(|\lambda| < \varrho(u)\) ;
\item
  The only eigenvectors of u in C are the multiples of \(x_0\) .
\end{enumerate}

Set \(\varrho = \varrho(u)\) . One can apply Th. 4 to \(u\) and Prop. 8 to \(u^k\) . There therefore exists an eigenvector \(x_0\) of \(u\) in C for the eigenvalue \(\varrho\) . Since \(0 < \varrho(u^k) = \varrho^k\) (formula (4) of I, p. 21), one has in particular \(\varrho > 0\) , and since \(x_0\) is an eigenvector of \(u^k\) for the eigenvalue \(\varrho(u^k)\) , we have \(x_0 \in \mathring{\mathrm{C}}\) .

One has \(\mathrm{Sp}(u_{(\mathbf{C})}^{k})=\mathrm{Sp}(u_{(\mathbf{C})})^{k}\) (Remark 4 of I, p. 2), hence every eigenvalue \(\lambda\in C\) of \(u_{(\mathbf{C})}\) such that \(|\lambda|=\varrho\) satisfies \(\lambda^{k}=\varrho^{k}\) , and since every eigenvector corresponding to \(\lambda\) is an eigenvector of \(u^{k}\) for \(\varrho^{k}\) , hence proportional to \(x_{0}\) , we have \(\lambda=\varrho\) .

If the spectral projector of u for the eigenvalue \(\varrho\) had rank at least 2, then so would that of \(u^{k}\) for the eigenvalue \(\varrho^{k}\) (cf. no. 2 of I, p. 129).

Finally, if \(x \in C\) is an eigenvector of u, it is also one for \(u^{k}\) , hence is proportional to \(x_{0}\) .

Let n be an integer \(\geqslant 1\) . The set \(R_{+}^{n}\) is a closed convex pointed salient cone in \(R^{n}\) , with nonempty interior equal to \((\mathbf{R}_{+}^{*})^{n}\) .

Let \(A = (a_{i,j})\) be a real matrix of type \((n,n)\) , and \(u_A\) the endomorphism \(x \mapsto Ax\) of \(\mathbf{R}^n\) . Let \(\varrho = \varrho(u_A)\) be its spectral radius.

Lemma 7. --- One has \(u_A(\mathbf{R}_+^n) \subset \mathbf{R}_+^n\) if and only if \(a_{i,j} \geqslant 0\) for all i and j, and \(u_A(\mathbf{R}_+^n - \{0\}) \subset (\mathbf{R}_+^*)^n\) if and only if \(a_{i,j} > 0\) for all i and j.

Let \((e_{1},\ldots,e_{n})\) be the canonical basis of \(R^{n}\) . The vectors \(e_{i}\) belong to \(R_{+}^{n}\) and \(u_{A}(e_{i})\in\mathbf{R}_{+}^{n}\) for all i (resp. \(u_{A}(e_{i})\in(\mathbf{R}_{+}^{*})^{n}\) for all i) if and only if \(a_{i,j}\geqslant0\) for all i and j (resp. \(a_{i,j}>0\) for all i and j). The result follows by linearity.

THEOREM 5 (Perron--Frobenius). --- a) If \(a_{i,j} \geqslant 0\) for all \(i\) and \(j\) in \(\{1,\ldots,n\}\) , then the real number \(\varrho\) is an eigenvalue of \(A\) ; b) If \(a_{i,j} > 0\) for all \(i\) and \(j\) in \(\{1,\ldots,n\}\) , then \(\varrho > 0\) and the primary space of \(u_A\) relative to \(\varrho\) , that is, the nullspace of \(A - \varrho 1_{\mathbf{R}^n}\) , has dimension 1 and is generated by a vector \(x_0 \in (\mathbf{R}_+^*)^n\) . There is no other eigenvalue of \(A\) admitting an eigenvector in \(\mathbf{R}_+^n\) , and all complex eigenvalues \(\lambda \neq \varrho\) of \(A\) satisfy \(|\lambda| < \varrho\) .

If \(\varrho = 0\) , assertion \(a\) ) is elementary because 0 is then an eigenvalue of \(A\) . If \(\varrho > 0\) , assertion \(a\) ) follows from Theorem 4 in view of Lemma 7. Assertion \(b\) ) is, for the same reason, a consequence of Proposition 8.

Remark. --- Suppose there exists an integer \(k \geqslant 1\) such that all coefficients of \(A^{k}\) are \textgreater0 (such a matrix is sometimes called primitive). Then, by the corollary of Prop. 8, the spectral radius \(\varrho\) is an eigenvalue of A, the primary space of \(u_{A}\) relative to \(\varrho\) has dimension 1, and is generated by a vector \(x_{0} \in (\mathbf{R}_{+}^{*})^{n}\) . Moreover, there is no other complex eigenvalue of A admitting an eigenvector in \(R_{+}^{n}\) and all complex eigenvalues \(\lambda \neq \varrho\) of A satisfy \(|\lambda| < \varrho\) .

